\documentclass{beamer}

\usepackage{graphicx}
\usepackage{amsmath}

\title{Immersed Boundary Method Discretization}
\author{Laura Miller}
\date{\today}

\begin{document}

\frame{\titlepage}


\begin{frame}{Discretization of the Immersed Boundary Method}
\small
The immersed boundary (IB) method discretizes a coupled
fluid–structure system with:
\begin{itemize}
    \item an \textbf{Eulerian grid} for the fluid variables
    \item a \textbf{Lagrangian mesh} for the immersed structure
    \item \textbf{regularized delta functions} to couple the two
\end{itemize}

\medskip
At each time step, the method alternates between:
\begin{enumerate}
    \item solving the fluid equations on a fixed grid,
    \item computing elastic forces on the structure,
    \item transferring forces and velocities between grids.
\end{enumerate}
\end{frame}

\begin{frame}{Eulerian Discretization: Fluid Variables}
\small
The fluid equations are discretized on a fixed Cartesian grid
\[
\mathbf{x}_{i,j} = (i h, j h)
\quad \text{(or } \mathbf{x}_{i,j,k} \text{ in 3D)}
\]


\medskip
\begin{itemize}
    \item Velocity: $\mathbf{u}_{i,j}^n \approx \mathbf{u}(\mathbf{x}_{i,j}, t^n)$
    \item Pressure: $p_{i,j}^n \approx p(\mathbf{x}_{i,j}, t^n)$
    \item Grid spacing: $h$
\end{itemize}

\medskip
We use $h$ and $\Delta x$ interchangeably to denote the uniform Eulerian grid spacing.

\medskip
Spatial derivatives are typically approximated using:
\begin{itemize}
    \item second-order finite differences,
    \item staggered (MAC) grids in many IB/IBAMR implementations.
\end{itemize}
\end{frame}

\begin{frame}{Lagrangian Discretization: Structure Variables}
\small
The immersed structure is discretized using material points
\[
\mathbf{X}_k^n \approx \boldsymbol{\chi}(\mathbf{q}_k, t^n),
\]
where $\mathbf{q}_k$ are reference coordinates.

\medskip
\begin{itemize}
    \item Lagrangian spacing: $\Delta s$ (1D), $\Delta A$ (2D), or $\Delta V$ (3D)
    \item Elastic force density:
    \[
    \mathbf{F}_k^n \approx \mathbf{F}(\mathbf{q}_k, t^n)
    \]
\end{itemize}

\medskip
Elastic forces are computed entirely in Lagrangian coordinates,
using springs, beams, shells, or continuum elasticity models.
\end{frame}

\begin{frame}{Force Spreading: Lagrangian to Eulerian}
\small
Elastic forces computed on the structure are spread to the fluid grid via
\[
\mathbf{f}_{i,j}^n
=
\sum_k
\mathbf{F}_k^n\,
\delta_h(\mathbf{x}_{i,j} - \mathbf{X}_k^n)\,
\Delta s
\]

\medskip
Here $\Delta s$ denotes the appropriate Lagrangian quadrature weight
(e.g.\ $\Delta s=\Delta q$ for a fiber, $\Delta A$ for a surface).


\medskip
\begin{itemize}
    \item $\delta_h$ is the regularized delta function
    \item Each Lagrangian force contributes to nearby fluid grid points
    \item Compact support ensures locality and efficiency
\end{itemize}

\medskip
This discretization approximates the continuous force-spreading integral
introduced earlier.
\end{frame}

\begin{frame}{Velocity Interpolation: Eulerian to Lagrangian}
\small
The structure moves with the local fluid velocity, interpolated as
\[
\mathbf{U}_k^n
=
\sum_{\mathbf{i}}
\mathbf{u}_{\mathbf{i}}^n\,
\delta_h(\mathbf{x}_{\mathbf{i}} - \mathbf{X}_k^n)\,
h^{d}
\]


\medskip
\begin{itemize}
    \item Uses the \emph{same} regularized delta function as force spreading
    \item Ensures consistency between force and velocity transfer
    \item Enforces the no-slip condition in a weak (integral) sense
\end{itemize}
\end{frame}

\begin{frame}{Updating the Structure}
\small
Once the velocity at each Lagrangian point is known,
the structure is advanced in time:
\[
\mathbf{X}_k^{n+1}
=
\mathbf{X}_k^n
+
\Delta t\, \mathbf{U}_k^n
\]

\medskip
\begin{itemize}
    \item Forward Euler shown for simplicity
    \item Higher-order and implicit schemes are also used in practice
    \item Time step restricted by elastic stiffness and fluid stability
\end{itemize}
\end{frame}

\begin{frame}{Immersed Boundary Algorithm (One Time Step)}
\small
Given $(\mathbf{u}^n, \mathbf{X}^n)$:

\begin{enumerate}
    \item Compute elastic forces $\mathbf{F}^n$ on the structure
    \item Spread forces to the fluid grid to obtain $\mathbf{f}^n$
    \item Solve the fluid equations for $\mathbf{u}^{n+1}$
    \item Interpolate velocity to obtain $\mathbf{U}^{n+1}$
    \item Update the structure configuration $\mathbf{X}^{n+1}$
\end{enumerate}

\medskip
This cycle is repeated for each time step.
\end{frame}

\begin{frame}{Time Step Restrictions and Stiffness}
\small

While conceptually simple, the Peskin (1996) discretization
imposes practical limitations on the time step.

\medskip
\textbf{Sources of time step restriction:}
\begin{itemize}
    \item Advection CFL condition (explicit convection)
    \item Strong stiffness from elastic forces
    \item Additional constraints depending on the time integrator
\end{itemize}


\medskip
For fiber bending forces:
\[
\mathbf{F}_{\text{bend}} \sim \frac{1}{(\Delta q)^4},
\]
leading to severe restrictions on $\Delta t$ when treated explicitly.

\medskip
This motivates later developments:
\begin{itemize}
    \item semi-implicit and fully implicit IB schemes
    \item IBFE and penalty formulations
    \item adaptive and multilevel solvers
\end{itemize}

\end{frame}

\begin{frame}{Codimension Perspective in the IB Method}
\small

The immersed boundary framework naturally handles structures of
different dimensions immersed in a higher-dimensional fluid.

\medskip
\begin{itemize}
    \item 1D fibers immersed in 2D or 3D fluids
    \item 2D surfaces immersed in 3D fluids
    \item (In some extensions) volumetric solids
\end{itemize}

\medskip
The same coupling formulas apply in all cases:
\[
\mathbf{f}(\mathbf{x})
=
\int_{\mathcal{U}}
\mathbf{F}(\mathbf{X})\,
\delta(\mathbf{x}-\boldsymbol{\chi}(\mathbf{X}))\,d\mathbf{X},
\]
with appropriate quadrature weights.

\medskip
This flexibility is one of the defining strengths of the IB approach.

\end{frame}



\begin{frame}{From Overview to a Concrete Implementation}
\small

So far, we have described the immersed boundary method at a general level:
\begin{itemize}
    \item Eulerian–Lagrangian coupling via integral transforms
    \item Regularized delta functions for interaction
    \item General algorithmic structure of IB methods
\end{itemize}

\medskip

We now turn to a \textbf{specific and influential implementation}:
\begin{itemize}
    \item the fiber-based immersed boundary method described in \textbf{Peskin (1996)}
\end{itemize}

\medskip

\textbf{Important caveats:}
\begin{itemize}
    \item This discretization reflects design choices made for efficiency and robustness at the time
    \item Modern IB frameworks (e.g.\ IBAMR, IBFE) generalize or modify several of these steps
    \item Our goal is understanding the \emph{logic} of the discretization, not prescribing a single best practice
\end{itemize}

\medskip

With these caveats in mind, we now walk through Peskin’s discretization step by step.

\end{frame}

\begin{frame}{Scope and Assumptions of Peskin (1996)}
\small

Before diving into the details, it is important to be clear about the
assumptions underlying the classical fiber-based immersed boundary method
of Peskin (1996).

\medskip
\textbf{Key assumptions:}
\begin{itemize}
    \item Incompressible, Newtonian fluid
    \item Uniform Cartesian Eulerian grid
    \item Periodic computational domain (FFT-based solvers)
    \item Codimension-1 structures (fibers) immersed in the fluid
    \item Elastic forces computed in Lagrangian coordinates
    \item Semi-implicit time stepping (viscosity implicit, coupling explicit)
\end{itemize}

\medskip
These assumptions were chosen to emphasize robustness, simplicity,
and computational efficiency.

\end{frame}

%=======================================================
% Discretization (slide 70 style)
%=======================================================
\begin{frame}{Discretization of Navier-Stokes}
\small
On a uniform Eulerian grid with spacing $\Delta x$ and timestep $\Delta t$, a common semi-implicit update is:
\begin{align*}
\rho\left(
\frac{\mathbf{u}^{n+1}-\mathbf{u}^{n}}{\Delta t}
+ \mathcal{S}_{\Delta x}(\mathbf{u}^{n})\,\mathbf{u}^{n}
\right)
- \mathbf{D}^{0} p^{n+1}
&=
\mu\sum_{\alpha=1}^{d} D^{+}_{\alpha}D^{-}_{\alpha}\mathbf{u}^{n+1}
+\mathbf{f}^{n},\\
\mathbf{D}^{0}\cdot \mathbf{u}^{n+1} &= 0,
\end{align*}
where $d\in\{2,3\}$ is the spatial dimension.

\medskip
\textbf{Key point:} The system is linear in the unknowns $\mathbf{u}^{n+1}$ and $p^{n+1}$ once $\mathbf{f}^n$ and the explicit convection term are known, so FFT-based solvers are natural on periodic boxes.
\end{frame}

%=======================================================
% Discretization details (slide 71 style)
%=======================================================
\begin{frame}{Finite difference operators used above}
\small
Let $\mathbf{e}_{\alpha}$ denote the Cartesian basis vectors ($\alpha=1,\dots,d$).

\medskip
\textbf{Centered difference gradient components:}
\[
\left(D_{\alpha}^{0}\phi\right)(\mathbf{x})
=
\frac{\phi(\mathbf{x}+\Delta x\,\mathbf{e}_{\alpha})-\phi(\mathbf{x}-\Delta x\,\mathbf{e}_{\alpha})}{2\Delta x},
\qquad
\mathbf{D}^{0}\phi = (D_{1}^{0}\phi,\dots,D_{d}^{0}\phi).
\]

\textbf{Forward/backward differences:}
\[
\left(D_{\alpha}^{+}\phi\right)(\mathbf{x})
=
\frac{\phi(\mathbf{x}+\Delta x\,\mathbf{e}_{\alpha})-\phi(\mathbf{x})}{\Delta x},
\qquad
\left(D_{\alpha}^{-}\phi\right)(\mathbf{x})
=
\frac{\phi(\mathbf{x})-\phi(\mathbf{x}-\Delta x\,\mathbf{e}_{\alpha})}{\Delta x}.
\]

\textbf{Skew-symmetric discrete convection (energy-friendly form):}
\[
\mathcal{S}_{\Delta x}(\mathbf{u})\,\phi
=
\frac{1}{2}\,\mathbf{u}\cdot \mathbf{D}^{0}\phi
+
\frac{1}{2}\,\mathbf{D}^{0}\cdot(\mathbf{u}\,\phi).
\]
\end{frame}

%=======================================================
% Boundary discretization (slide 72 style)
%=======================================================
\begin{frame}{Fiber/Boundary discretization (1D structure in $d$D fluid)}
\small
Discretize the immersed fiber by Lagrangian nodes $\mathbf{X}^{n}_{m}\approx \mathbf{X}(q_m,t^n)$
with spacing $\Delta q$ (using $q$ as the Lagrangian parameter to avoid confusion with meters).

\medskip
Common force components (additive):
\[
\mathbf{F}^{n}_{m}
=
\mathbf{F}^{n}_{\text{targ},m}
+
\mathbf{F}^{n}_{\text{str},m}
+
\mathbf{F}^{n}_{\text{bend},m}.
\]

\textbf{Target-point (penalty) force:}
\[
\mathbf{F}^{n}_{\text{targ},m}
=
k_{\text{targ}}\left(\mathbf{Y}^{n}_{m}-\mathbf{X}^{n}_{m}\right),
\]
where $\mathbf{Y}^{n}_{m}$ is a prescribed target location.

\medskip
\textbf{Stretching force (Hookean tension along the fiber):}
Define the centered difference
\[
\left(D_{q}\mathbf{X}^{n}\right)_m
=
\frac{\mathbf{X}^{n}_{m+1}-\mathbf{X}^{n}_{m-1}}{2\Delta q},
\qquad
\boldsymbol{\tau}^{n}_m
=
\frac{(D_{q}\mathbf{X}^{n})_m}{\left\lVert (D_{q}\mathbf{X}^{n})_m\right\rVert}.
\]
A common discrete stretching force is
\[
\mathbf{F}^{n}_{\text{str},m}
=
k_{\text{str}}\,D_{q}\!\left(\left(\left\lVert D_{q}\mathbf{X}^{n}\right\rVert-1\right)\boldsymbol{\tau}^{n}\right)_m.
\]


\end{frame}

\begin{frame}{Fiber/Boundary discretization (1D structure in $d$D fluid)}
\textbf{Bending (beam) force (discrete 4th derivative):}
\[
\mathbf{F}^{n}_{\text{bend},m}
=
k_{\text{bend}}\,
\frac{\mathbf{X}^{n}_{m-2}-4\mathbf{X}^{n}_{m-1}+6\mathbf{X}^{n}_{m}-4\mathbf{X}^{n}_{m+1}+\mathbf{X}^{n}_{m+2}}{(\Delta q)^4}.
\]
    
\end{frame}
%=======================================================
% Spreading + interpolation + update (slide 73 style)
%=======================================================
\begin{frame}{Force spreading and velocity interpolation (discrete)}
\small
Let Eulerian grid points be $\mathbf{x}_{\mathbf{i}}$ (multi-index $\mathbf{i}$ in $d$D).

\medskip
\textbf{Spread Lagrangian force to the Eulerian grid:}
\[
\mathbf{f}^{n}(\mathbf{x}_{\mathbf{i}})
=
\sum_{m}\mathbf{F}^{n}_{m}\;
\delta_{\Delta x}\!\left(\mathbf{x}_{\mathbf{i}}-\mathbf{X}^{n}_{m}\right)\,\Delta q.
\]

\textbf{Interpolate Eulerian velocity to the fiber:}
\[
\mathbf{U}^{n+1}_{m}
=
\sum_{\mathbf{i}}\mathbf{u}^{n+1}(\mathbf{x}_{\mathbf{i}})\;
\delta_{\Delta x}\!\left(\mathbf{x}_{\mathbf{i}}-\mathbf{X}^{n}_{m}\right)\,(\Delta x)^{d}.
\]

\textbf{No-slip update of the structure:}
\[
\frac{\mathbf{X}^{n+1}_{m}-\mathbf{X}^{n}_{m}}{\Delta t}
=
\mathbf{U}^{n+1}_{m}
\qquad\Longleftrightarrow\qquad
\mathbf{X}^{n+1}_{m}=\mathbf{X}^{n}_{m}+\Delta t\,\mathbf{U}^{n+1}_{m}.
\]

\medskip
Here $\delta_{\Delta x}$ is a regularized approximation of the $d$-D Dirac delta.
\end{frame}

\begin{frame}{Weak Enforcement of the No-Slip Condition}
\small

In the immersed boundary method, the no-slip condition is not enforced
pointwise at the structure.

\medskip
Instead:
\begin{itemize}
    \item Fluid velocity is interpolated to Lagrangian points using $\delta_h$
    \item The structure is advanced using this interpolated velocity
    \item Forces are spread back to the fluid using the same kernel
\end{itemize}

\medskip
As a result, no-slip is enforced \textbf{in an averaged (integral) sense}:
\[
\frac{d\mathbf{X}}{dt}(\mathbf{q},t)
\approx
\int \mathbf{u}(\mathbf{x},t)\,
\delta_h(\mathbf{x}-\boldsymbol{\chi}(\mathbf{q},t))\,d\mathbf{x}.
\]

\medskip
Small grid-scale slip may occur, but vanishes as $h \to 0$.

\end{frame}

%=======================================================
% Discrete delta function (slide 74 style) -- cosine kernel
%=======================================================
\begin{frame}{A common discrete delta: tensor-product cosine kernel}
\small
A standard choice is a tensor-product kernel:
\[
\delta_{\Delta x}(\mathbf{x})
=
\frac{1}{(\Delta x)^{d}}
\prod_{\alpha=1}^{d}\phi\!\left(\frac{x_{\alpha}}{\Delta x}\right),
\qquad d\in\{2,3\}.
\]

\medskip
\textbf{Cosine (4-point support) 1D kernel:}
\[
\phi(s)=
\begin{cases}
\frac{1}{4}\left(1+\cos\!\left(\frac{\pi s}{2}\right)\right), & |s|\le 2,\\[6pt]
0, & |s|>2.
\end{cases}
\]

\medskip
\textbf{Useful properties (designed for IB coupling):}
\begin{itemize}
\item Symmetric: $\phi(s)=\phi(-s)$.
\item Compact support: only nodes within $2\Delta x$ contribute in each coordinate.
\item Partition of unity (discrete “mass conservation”):
\[
\sum_{j\in\mathbb{Z}}\phi(s-j)=1.
\]
\end{itemize}
\end{frame}

\begin{frame}{From Peskin (1996) to Modern IB Methods}
\small

The discretization we have examined is foundational, but not the end
of the story.

\medskip
\begin{center}
\begin{tabular}{l|c|c}
\textbf{Feature} & \textbf{Peskin (1996)} & \textbf{Modern IB (e.g.\ IBAMR)} \\
\hline
Structure type & Fibers & Fibers, surfaces, solids \\
Elasticity & Discrete springs/beams & FE-based continuum models \\
Coupling & Mostly explicit & Semi-/fully implicit \\
Solvers & FFT, periodic & Multilevel, adaptive \\
\end{tabular}
\end{center}

\medskip
Understanding the classical method provides essential insight into
why modern IB methods are designed the way they are.

\end{frame}

\begin{frame}{Roadmap: From Classical IB to Modern Methods}
\small

The discretization presented here follows Peskin (1996) and provides a
clear conceptual foundation for immersed boundary methods.

\medskip
\textbf{What this formulation does well:}
\begin{itemize}
    \item Robust Eulerian–Lagrangian coupling
    \item Simple handling of moving, deformable structures
    \item Clear physical interpretation of forces and motion
\end{itemize}

\medskip
\textbf{Key limitations of the classical approach:}
\begin{itemize}
    \item Severe time step restrictions for stiff structures
    \item Limited accuracy near immersed boundaries
    \item Periodic-domain and uniform-grid assumptions
\end{itemize}

\medskip
\textbf{Natural extensions developed later:}
\begin{itemize}
    \item Implicit and semi-implicit coupling strategies
    \item Finite-element structure models (IBFE)
    \item Adaptive mesh refinement and multilevel solvers
\end{itemize}

\end{frame}

\end{document}